\documentclass[12pt]{article}
\usepackage{amsmath}
\usepackage{mathrsfs}
\usepackage{eucal}
\usepackage{amssymb}
\usepackage[french]{babel}
\usepackage[latin1]{inputenc}
\usepackage[T1]{fontenc}
\usepackage[all]{xy}
%\usepackage{graphicx}
%\usepackage{epsf}

\begin{document}

\hfuzz 8 pt
\parindent = 0 cm

\newcommand{\card}[1]{\overline{\overline{#1}}}

\centerline{ \textsc{ \large Universit� de Sherbrooke}}
\centerline{\textsc{ \normalsize D�partement de math�matiques}}
%\centerline{\textsc{ \normalsize Automne 2005}}
\bigskip

\centerline{\large\bf MAT501 -- Fondements et histoire des
math�matiques}
\bigskip

\centerline{\Large\bf Laboratoire 1}
\bigskip

Construire � la r�gle et au compas en justifiant chaque �tape �
l'aide des axiomes de Hilbert :

\begin{description}

\item[1)] La m�diatrice d'un segment donn� ;

\item[2)] Le milieu d'un segment ;

\item [3)] La bissectrice d'un angle donn� ;

\item[4)] La perpendiculaire par un point $P$ � une droite $d$
avec $P \in d$ ;

\item[5)] La perpendiculaire par un point $P$ � une droite $d$
avec $P \notin d$ ;

\item[6)] La parall�le par un point $P$ � une droite $d$ avec
$P\notin d$ ;

\item[7)] Un angle de $\pi/3$, de $\pi/4$ et de $\pi/6$ ;

\item[8)] Un triangle �quilat�ral sur un segment donn� ;

\item[9)] Un carr� sur un segment donn� ;

\item[10)] Une cercle par trois points non colin�aires donn�s ;

\item[11)] Le cercle inscrit d'un triangle donn� ;

\item[12)] Le tiers d'un segment donn� ;

\item[13)] Le tiers d'un angle donn� ;

\item[14)] Un hexagone r�gulier sur un segment donn� ;

\item[15)] Un octogone r�gulier sur un segment donn� ;

\item[16)] � partir d'un segment $AB$ de longueur 1 un segment de
longueur :

  \begin{description}

  \item [a)] $m/n$ avec $n,m \in \mathbb{N}$ et $n \neq 0$ ;

  \item[b) ] $\sqrt{2}$ ;

  \item[c)] $\sqrt{3}$ ;

  \item[d)] $\sqrt{m}$ avec $m \in \mathbb{N}$.

  \end{description}

\end{description}

\end{document}
